% Source: https://artofproblemsolving.com/community/c3h2704867_2014_amc_10b_24
% Stable identity: aops-topic-2704867-q01
Up to rotation and reflection there are $\dfrac{4!}{2}=12$ arrangements. In every arrangement the single numbers give $1,\dots,5$, the whole circle gives $15$, and the complements of single numbers give $10,\dots,14$. A consecutive block summing to $n$ leaves a consecutive block summing to $15-n$, so only $6$ and $7$ need checking (they give $9$ and $8$).

\emph{Missing $6$.} The ways to make $6$ are $\{1,5\}$, $\{2,4\}$ and $\{1,2,3\}$, so $1,5$ must not be adjacent, $2,4$ must not be adjacent, and $1,2,3$ must not be consecutive. A short check of the placements satisfying these three restrictions leaves only the circle $1,2,5,3,4$ (up to symmetry).

\emph{Missing $7$.} The ways are $\{3,4\}$, $\{2,5\}$ and $\{1,2,4\}$; the same reasoning forces $1,3,2,4,5$.

These two arrangements are different, so there are $\boxed{\textbf{(B) }2}$ bad arrangements.

(Motivation: the complement trick reduces $15$ targets to $2$. For each remaining target, list its few decompositions; "bad" means none of them is consecutive, and these adjacency restrictions are strong enough to pin the arrangement down.)
